\chapter{A Tour of Zeno}\label{ch-tour}

This chapter runs through the whole language quickly, one small
program at a time, so that you know what is coming and where it is
explained. Nothing here has to be understood yet. Each section names
the chapter that goes through it properly.

\section{Values that change, and values that do not}

\begin{zeno}
let name = "Zeno"         # let: set once, never changes
var count = 3             # var: can change later
count = count + 1         # plain = changes a var
const LIMIT = 100         # const: fixed when the program is built
print(name, count, LIMIT)
\end{zeno}

A name is declared once, with \verb|let|, \verb|var| or \verb|const|,
and from then on plain \verb|=| changes it. A misspelt name is
therefore a mistake, and never a new variable.
(\hyperref[ch-values]{Values and Variables}.)

\section{Lists, and a comprehension}

\begin{zeno}
let readings = [21.5, 22.0, -99.0, 21.0, 150.0]
let kept = [r for r in readings if -40.0 <= r <= 85.0]
print(kept, sum(kept) / len(kept))
\end{zeno}

Comparisons chain as they do in mathematics: \verb|-40.0 <= r <= 85.0|
means what it looks like. (\hyperref[ch-lists]{Lists},
\hyperref[ch-numbers]{Numbers and Operators}.)

\section{Functions, with types or without}

\begin{zeno}
function describe(t: float) -> str {
    if t < 18.0 {
        return "cold"
    } else if t < 24.0 {
        return "comfortable"
    }
    return "warm"
}

function average(values) {
    return sum(values) / len(values)
}

for t in [12.0, 21.5, 30.0] {
    print(t, describe(t))
}
print(average([1, 2, 3, 4]))
\end{zeno}

Types are optional. Where they are written, they are checked before
the program runs. Where they are left out, Zeno works them out from
how the function is used. (\hyperref[ch-functions]{Functions},
\hyperref[ch-types]{Types, Written and Worked Out}.)

\section{Classes}

\begin{zeno}
import math

class Point {
    x: float
    y: float

    function distance_to(self, other: Point) -> float {
        return math.sqrt((self.x - other.x) ** 2 + (self.y - other.y) ** 2)
    }
}

let origin = Point(0.0, 0.0)
let here = Point(x=3.0, y=4.0)
print(origin, here, origin.distance_to(here))
\end{zeno}

A class lists its fields, and gets a way to build one for free. It has
no parent class: what classes share is said with an interface.
(\hyperref[ch-classes]{Classes}, \hyperref[ch-interfaces]{Interfaces}.)

\section{Enums that carry values, and match}

\begin{zeno}
enum Reading {
    Value(celsius: float)
    Fault(code: int)
    Missing
}

function show(r: Reading) -> str {
    match r {
        case Value(c) { return f"{c:.1f} degrees" }
        case Fault(code) { return f"fault {code}" }
        case Missing { return "no reading" }
    }
}

for r in [Value(21.5), Fault(7), Missing] {
    print(show(r))
}
\end{zeno}

Leave a \verb|case| out and Zeno names it before the program runs.
(\hyperref[ch-enums]{Enums and match}.)

\section{No surprise none}

\begin{zeno}[fails=yes]
function find_user(id: int) -> str or none {
    if id == 1 {
        return "ada"
    }
    return none
}

let user = find_user(2)
print(user.upper())
\end{zeno}

A value can be \verb|none| only if its type says \verb|or none|, and
then Zeno will not let it be used until the program has asked:

\begin{zeno}
function find_user(id: int) -> str or none {
    if id == 1 {
        return "ada"
    }
    return none
}

for id in [1, 2] {
    let user = find_user(id)
    if user is none {
        print(id, "nobody")
    } else {
        print(id, user.upper())
    }
}
\end{zeno}

(\hyperref[ch-none]{Values That May Be Missing}.)

\section{Errors that are checked}

\begin{zeno}
function parse_port(text: str) -> int raises ValueError {
    let port = int(text)
    if not 1 <= port <= 65535 {
        raise ValueError(f"{port} is not a port")
    }
    return port
}

for text in ["8080", "99999", "http"] {
    try {
        print("port", parse_port(text))
    } except ValueError as err {
        print("refused:", err)
    }
}
\end{zeno}

A function says what it \verb|raises|, and Zeno checks that every error
is either caught or passed on on purpose. (\hyperref[ch-errors]{Errors}.)

\section{Functions are values}

\begin{zeno}
function counter() {
    var count = 0
    return function() {
        count += 1
        return count
    }
}

let next_id = counter()
print(next_id(), next_id(), next_id())
print([1, 2, 3].map(function(n: int) -> int { return n * 2 }))
\end{zeno}

(\hyperref[ch-function-values]{Functions as Values}.)

\section{Dicts}

\begin{zeno}
let ages = ["Ada": 36, "Alan": 41]
ages["Grace"] = 85
for name, age in ages {
    print(name, age)
}
\end{zeno}

(\hyperref[ch-dicts]{Dicts}.)

\section{Numbers with units, and numbers with a size}

\begin{zeno}
let delay = 1500 ms
print(delay, delay.s, 1 s + 250 ms)

let frame: u8 = 200
print(u8.wrap(frame + 100), frame & 0x0F, frame >> 4)
\end{zeno}

A number can carry a unit, and units that measure the same thing work
together. A whole number can have a size, and a number that does not
fit is an error unless you ask for it to wrap.
(\hyperref[ch-units]{Units}, \hyperref[ch-sized-numbers]{Sized Numbers
and Bits}.)

\section{Generics}

\begin{zeno}
function largest[T](items: list[T]) -> T {
    var best = items[0]
    for item in items {
        if item > best {
            best = item
        }
    }
    return best
}

print(largest([3, 9, 4]), largest(["pear", "fig"]), largest([3.3 V, 5 V]))
\end{zeno}

Written once, checked for each type it is used with, and compiled as
if it had been written by hand for that type.
(\hyperref[ch-generics]{Generics}.)

\section{And on a board}

\begin{zeno}[file=blink.zeno, cmd=zeno check blink.zeno]
from board import led
from time import sleep

while true {
    led.toggle()
    sleep(500 ms)
}
\end{zeno}

On a computer the board is simulated, and \verb|zeno run| prints each
change of a pin. The same program builds for firmware with no garbage
collector, no heap and no \verb|setjmp|. (\hyperref[ch-board]{Pins,
Time and the Board}, \hyperref[ch-firmware]{Firmware Builds},
\hyperref[ch-no-collector]{Building with No Collector}.)

\begin{tryit}
\begin{enumerate}
\item Take one of the programs in this chapter, break it on purpose,
  and read the message.
\item In the enum program, add a fourth case \verb|Calibrating| to
  \verb|Reading| and run it before changing \verb|show|.
\item In the error program, take \verb|raises ValueError| off
  \verb|parse_port| and run \verb|zeno check --strict| on the file.
\end{enumerate}
\end{tryit}
